\section{Datasets and Experimental Methodology}
\label{sec:methodology}

Each run trains and evaluates one model with a specified objective. A
constrained run is paired with a baseline using the same model recipe and
labelled training examples, called the support set. They also share the data
partition and the random settings, defined below. The medical four-model
studies match the number of floating-point operations (FLOPs) used during
optimization; one GFLOP is one billion such operations. The geological
four-model study caps training plus calibration. MedSAM3 and SAM3 instead
match adapter updates within their respective schedules. These task and
resource definitions delimit the comparisons that can be made.

\subsection{Design and Units of Comparison}
\label{sec:experimental-structure}

A fold identifies a partition into training and validation examples;
validation examples are held out from parameter fitting. A seed initializes
the pseudorandom generator controlling training randomness. Repeating seeds
within a fold measures sensitivity to initialization and data presentation;
changing folds also changes the support and evaluation images. These
repetitions provide evidence about consistency under the available data,
while sharing parts of the same experimental setting.

Table~\ref{tab:experimental-design} records the six studies, including two
geological and four medical comparisons. A condition is a particular model,
dataset and objective. Each condition has nine fold/seed cells, one run for
each combination of three folds and three seeds. A checkpoint is a saved
model state. Evaluation uses the final state allowed by the budget in the
capped experiments, the final 200-update MedSAM3 adapter and the final
25-update SAM3 adapter, without validation-selected checkpoint replacement
or early stopping.

Throughout the comparisons, \emph{Baseline} denotes the supervised training
objective configured for the particular model and dataset. \emph{Baseline}
$+B$, for example, retains that objective and adds boundary-band
supervision; the same convention applies to combinations of spatial terms.
The baseline may itself combine several losses, as specified in the
model-specific protocols below. In tables and charts, $B$, $E$, $S$ and $D$
abbreviate boundary bands, edges, signed distance and depth supervision,
respectively.

\begin{table}[htbp]
\centering\small\setstretch{1.05}
\begin{tabular}{@{}P{2.5cm}P{3.0cm}P{3.2cm}P{5.7cm}@{}}
\hline
Study & Support per fold & Folds; seeds & Conditions and budget \\
\hline
GeoFWI3D, from scratch & 19 scenes & 0,1,2; 0,1,2 & 4 models $\times$ 14 objectives; 504 runs; $182{,}178.13$ GFLOPs/run including calibration. \\[.5em]
MSD, from scratch & 10 crops & 0,1,2; 0,2,20 & 4 models $\times$ 14 objectives; 504 runs; $31{,}976.45$ GFLOPs/run. \\[.5em]
ADNI, from scratch & 10 crops, 5 subjects & 0,1,2; 0,1,4 & 4 models $\times$ 14 objectives; 504 runs; same FLOP cap. \\[.5em]
GeoFWI3D, SAM3 & 5 scenes & 0,1,3; 0,1,2 & 14 objectives; 126 runs; 25 updates/run. \\[.5em]
MSD, MedSAM3 & 5 crops & 0,1,3; 17,83,191 & 14 objectives; 126 runs; 200 updates/run. \\[.5em]
ADNI, MedSAM3 & 2 crops, 2 subjects & 0,1,3; 0,1,2 & 14 objectives; 126 runs; 200 updates/run. \\
\hline
\end{tabular}
\caption{Completed study design: 1,890 objective runs in 210 conditions,
comprising 630 geological and 1,260 medical runs.
Each MedSAM3 campaign covers fourteen objectives, counting the baseline
objective once. Repeated matched controls (baseline runs) support paired gains but do
not enlarge the objective inventory.
Geological recipes and strengths were tuned during development;
MedSAM3 fold selection is exploratory.
All scores use matched final-budget baselines.}
\label{tab:experimental-design}
\end{table}

The medical conventional seed sets were chosen after inspecting
preliminary outcomes, so these comparisons are exploratory.
The numeric seed values identify pseudorandom streams; equal numbers
across different models or datasets would not make their initial states
identical. Each study retains its recorded seed inventory for reproducibility.
Within each dataset, all four architectures and all fourteen objectives
share the same seed set,
preserving matched comparisons between each constrained objective and
its baseline. Each objective includes all nine combinations of three
folds and three seeds, with equal weight for the three folds.

An annotated \emph{volume} is the stored example. In MSD, this is a
hippocampal crop, so a ten-volume or five-volume support set is an
annotation budget in crops. The available identifiers do not establish
how those crops map to participants. ADNI retains subject identifiers
and allows subject grouping to be checked explicitly.
Table~\ref{tab:dataset-task-overview} summarizes the input types and label
vocabularies. Each dataset section below illustrates these with an image,
its reference mask and an overlay of the two. The figures show slices
through the native three-dimensional volumes, before model-specific
padding or cropping. Only MSD includes an internal anterior/posterior division.
The experiments evaluate segmentation, with no downstream diagnostic task.
Within the chosen viewing direction, each example shows the slice with
the largest labelled area. The colours are transparent in the overlays
so that the underlying image remains visible. Image and label panels use
the same view and preserve the proportions of the stored grid.

\begin{table}[htbp]
\centering
\small
\setstretch{1.05}
\renewcommand{\arraystretch}{1.2}
\begin{tabular}{P{2.65cm}P{3.15cm}P{3.65cm}P{4.1cm}}
\hline
Dataset & Model input & Reference classes & Role in the investigation \\
\hline
GeoFWI3D & Synthetic 3D seismic reflectivity &
$0$: background; $1$: salt &
Nonmedical four-model comparison and SAM3 adaptation. \\[0.4em]
MSD hippocampus & Single-channel MRI crop &
$0$: background; $1$: anterior; $2$: posterior &
Main medical comparison, including the FLOP-capped study and MedSAM3 adaptation. \\[0.4em]
ADNI hippocampus & Single-channel MRI; separate left/right crops &
$0$: background; $1$: hippocampus &
Second medical dataset for whole-hippocampus segmentation. \\
\hline
\end{tabular}
\caption{Dataset tasks and class encodings in the local preparations.
GeoFWI3D uses reflectivity as input; velocity and relative geological time
are used only to construct its reference salt mask. ADNI laterality
identifies separate crops and does not create separate foreground classes.
Experimental subsets and partitioning are specified in this chapter.}
\label{tab:dataset-task-overview}
\end{table}


\FloatBarrier
\subsection{GeoFWI3D: Salt Segmentation}
\label{sec:geo-methods}


The geological collection is drawn from archive 0 of GeoFWI3D
\cite{swaminadhan2026geofwi3d}, containing 1,000 synthetic scenes on
$96\times96\times96$ grids. The experiment retains 479 scenes with
salt: 249 without faults and 230 with faults. This selection defines
a delineation task conditional on salt being present. Entirely
salt-free scenes are outside its evaluation population.

Each example pairs a reflectivity volume with a binary annotation that
identifies salt. Figure~\ref{fig:geo-data-examples} shows one cross-section:
the image contains the measurements available to the model, and the blue
region marks the salt body it is trained to identify.

\begin{figure}[htbp]
\centering
% Native data; generator: figures/source/dataset-methods-examples.py.
% Provenance: reports/thesis_dataset_methods_examples_20261103/provenance.json.
\includegraphics[width=0.98\linewidth]{figures/compiled/geo-data-examples.pdf}
\caption{One GeoFWI3D example: reflectivity image, reference salt mask and
their overlay on the same slice. Blue denotes salt; white in the mask
denotes background. The slice is taken from a $96\times96\times96$
volume, with geological depth increasing downward. Reflectivity is
displayed over $[-1,1]$. The image illustrates the input and training
annotation, not a model prediction.}
\label{fig:geo-data-examples}
\end{figure}

For preprocessing, each reflectivity volume is divided by
the 99.5th percentile of its absolute intensities and clipped to
$[-1,1]$, using only image values. Velocity and relative geological time (RGT), a coordinate describing
the ordering of geological layers, are used to build the annotation. Inspection showed that zero
relative geological time can also mark null artifacts, so the reference
requires agreement with the scene's recorded salt velocity. Write
$V_{P,i}(v)$ for the velocity at voxel $v$ and $V_{\mathrm{salt},i}$
for the recorded salt value. The latter is converted to the same
32-bit floating-point format, denoted float32, as the velocity volume:
\begin{equation}
 y_i(v)=\mathbf{1}\{\mathrm{RGT}_i(v)=0\}
        \mathbf{1}\{V_{P,i}(v)=\operatorname{float32}(V_{\mathrm{salt},i})\}.
\label{eq:geo-reference-construction}
\end{equation}
The equality uses the stored simulation value. All physically valid
salt components are retained, and checks reject physical salt outside
the implicit mask. This correction matters for boundary supervision:
an artifact in the source field would otherwise be treated as a
geological interface.


The completed geological comparison evaluates nnU-Net, SwinUNETR,
UNETR and the compact 3D U-Net under the same fourteen-objective
inventory used in the medical studies. Folds 0, 1 and 2 each contain
19 labelled training scenes and 96 validation scenes on the native
$96^3$ grid. The support set is the fixed approximately 5\% subset
of the eligible training pool. Training and validation are disjoint
within each fold, and the three validation sets are pairwise disjoint.
The support sets contain 54 distinct scenes across their 57 entries.
Seeds 0, 1 and 2 are used in every condition, giving 504 cells.
All nine final-budget outcomes enter each condition mean.

The final matrix combines 36 matched frozen baselines with 468
constrained results. 

\subsubsection{Frozen Model Recipes and Spatial Geometry}

The model recipes use one volume per batch and mixed numerical precision:
selected operations use the 16-bit bfloat16 format, while parameters remain
in float32. Automatic choice of operation precision is called autocast.
Computation uses NVIDIA's CUDA platform on graphics processing units (GPUs).
nnU-Net uses its native Dice-plus-cross-entropy objective and deep
supervision. SwinUNETR and the compact U-Net use Dice with background
included. Geological UNETR uses foreground Dice plus cross-entropy, squared
prediction terms in Dice, dropout 0.1 and learning rate $6\times10^{-4}$.
Dropout randomly suppresses a fraction of intermediate activations during
training; the learning rate controls the size of parameter updates.

In preliminary experiments, the initial UNETR recipe consistently struggled
to achieve satisfactory validation performance across the seeds tested
within the fixed FLOP budget. Two modifications were introduced during
development to improve performance on unseen GeoFWI3D scenes: horizontal
data augmentation, which creates transformed presentations of training
examples, and exclusion of the observation frame from the spatial
constraints. Both address properties of this geological task and were
retained in its final UNETR training recipe.

The augmentation policy selects among eight transformations in the
horizontal (XY) plane: rotations of $0^\circ$, $90^\circ$, $180^\circ$
or $270^\circ$, each with or without a reflection. Each transformation
is applied jointly to the image and its reference mask, preserving the
depth (Z) axis and their alignment. These operations rearrange voxels
without interpolation or intensity changes and are applied only during
training. The policy is shared by the UNETR baseline and all its
constrained objectives.

The second modification excludes the outer two voxel layers from the
band and edge supports. A salt body can continue beyond the observed
cube, whereas zero-padded erosion can treat the cube face as a region
boundary. Exclusion is intended to concentrate auxiliary supervision
on interfaces observed within the volume, reducing its dependence on
where the observation window ends. It changes the support of the
constraints while preserving the full input, prediction grid and
baseline segmentation loss. Its development comparison is reported in
Appendix~\ref{sec:geo-development}; the modification is a task-specific
loss-design choice rather than a requirement of the UNETR architecture.

All four models use two-step, six-neighbour bands on binary salt
membership. The other three geological recipes retain zero-padded
morphology at the cube faces. All four use stabilized band denominators
and require both sides to be nonempty.
Appendix~\ref{sec:geo-boundaries} specifies the reductions. The final
comparison therefore evaluates complete model recipes, including these
choices, rather than varying architecture alone.

Signed-distance supervision uses one salt field in voxel units, clipped
at eight voxels and normalized by eight. Depth supervision uses the
source-array geometry: its plane normal is source axis 0 projected
orthogonally to the salt mask's first principal axis. If the projected
norm is below 0.1, the least-parallel source axis is substituted, with
axis-index tie breaking. Depth strata are two voxels wide, the region
margin is eight voxels, and sparse strata use a threshold of 64 voxels.
These widths match those used for ADNI and, in voxel units, for MSD's
$1\,\mathrm{mm}$ grid. Their numerical agreement supports comparison of
the constraint settings without assuming equal physical scales across
applications.
UNETR's coordinate mapping is restored after the horizontal transformations before
constructing depth weights. This is an annotation-derived coordinate;
it is not a measurement of geological age or a guaranteed vertical
stratigraphic ordering.

\subsubsection{Calibration, Scheduling and Compute Allowance}

The geological study uses a larger computation allowance because its
inputs are larger. The reference is the original measured cost of 152
nnU-Net updates on the $96^3$ grid, approximately 182,178 GFLOPs. This
corresponds to eight passes through the 19 support scenes under the
reference profile. The final recipes retain that numerical ceiling and
charge both calibration and optimization within it. Evaluation and recipe
development remain outside the budget. Medical and geological results
therefore do not constitute an equal-compute comparison across datasets.

Calibration measures gradients on all 19 support scenes. Its strengths
differ across architectures, and the nnU-Net and UNETR settings follow
exploratory comparisons. Auxiliary terms enter through fixed ramps, during
which their weights increase gradually from zero to their calibrated values.
Some signed-distance configurations start later. Each constrained run is
compared with the baseline of the same frozen model recipe. The final
admissible checkpoint supplies the score.
Appendix~\ref{sec:geological-budget-details} records the exact ceiling,
strengths and schedules, while Table~\ref{tab:cap-training-exposure} gives
the actual training exposure.

\subsection{MSD: Medical Segmentation Decathlon}
\label{sec:msd-data-methods}


The MSD collection contains 263 labelled MRI crops with background,
anterior and posterior hippocampus classes \cite{simpson2019dataset}.
Figure~\ref{fig:msd-data-examples} shows one crop and its annotation.
Together, the blue and orange regions form the hippocampus; their shared
interface is the internal division that the model must recover in addition
to the outer boundary.

\begin{figure}[htbp]
\centering
% Native data; generator: figures/source/dataset-methods-examples.py.
% Provenance: reports/thesis_dataset_methods_examples_20261103/provenance.json.
\includegraphics[width=0.98\linewidth]{figures/compiled/msd-data-examples.pdf}
\caption{One MSD example: an MRI slice, its reference labels and their
overlay. Blue marks the anterior hippocampus and orange the posterior
hippocampus; white in the mask denotes background. All three panels show
the same slice of the native $35\times51\times35$ crop. MRI contrast is
adjusted for display using the volume's first and ninety-ninth intensity
percentiles. The coloured regions are reference annotations, not model
predictions.}
\label{fig:msd-data-examples}
\end{figure}

For the capped experiments, each image is z-score normalized: the mean of
its nonzero intensities is subtracted and the result divided by their
standard deviation. The image and its labels are then padded (extended with
zeros) and centrally cropped to $64^3$. This preserves voxel scale rather
than resizing the anatomy. Validation uses this prepared grid. Each of folds
0, 1 and 2 supplies ten support crops and 52 validation crops, shared across
models and objectives. Seeds are 0, 2 and 20. Validation identifiers are
disjoint across the three folds, but the available crop identifiers do not
establish participant-level independence. MedSAM3 uses its separate
native-grid preparation and manifests, files listing the cases assigned to
each partition, in Section~\ref{sec:medsam-methods}.

\subsubsection{Testing Proposed Anatomical Rules}
\label{sec:annotation-compatibility}

The MSD data descriptor \cite{simpson2019dataset} cites the manual
annotation protocols of Pruessner \emph{et al.}\
\cite{pruessner2000volumetry} and Woolard and Heckers
\cite{woolard2012asymmetry}. It identifies the last hippocampal-head slice
as the coronal slice containing the uncal apex, the landmark used to mark
the end of the head; body and tail form the posterior label. Coronal planes
divide the anatomy into front and back portions; sagittal planes divide left
and right. Turning this description into an automatic constraint requires a
detectable landmark, a known orientation and a relation that remains
meaningful on the stored grid. The annotation audit tests these operational
assumptions before imposing them during learning.

Two analyses address the proposed A/P transition rules. A planar analysis tests
whether one coronal separation reproduces the stored labels in all
260 crops. Coronal-profile and image-cue analyses examine 52 original
validation cases for a reliable way to locate the transition.

The image-cue investigation includes automated and visual searches but
has no independently validated expert localization of the uncal apex.
The audit can therefore establish whether a tested numerical proxy fits
the available labels. It cannot determine whether the original annotators
misapplied the anatomical protocol. Appendix~\ref{sec:annotation-results}
reports the findings within the limitations and explains their
consequences for constraint design.


\subsection{ADNI: Alzheimer's Disease Neuroimaging Initiative}
\label{sec:adni-methods}

The ADNI dataset used here combines the final released hippocampal labels
and oriented images from the European Alzheimer's Disease Consortium
(EADC)--ADNI harmonized-protocol resource \cite{harpReleasedLabels}. Image
and label pairs are identified through subject and image metadata. Because
the release notes a possible alignment issue in converted NIfTI images, the
builder uses the original MINC-format image arrays to export paired crops.

The preparation contains 135 subject identifiers and 135 image
identifiers, each contributing left and right hippocampal crops.
The resulting 270 volumes have size $32\times48\times40$ and binary
labels: zero for background and one for hippocampus. Laterality
identifies the crop; it does not define a further segmentation class.
Figure~\ref{fig:adni-data-examples} shows paired left and right crops,
their binary reference masks and annotation overlays. Both sides use
the same foreground class, without the anterior/posterior subdivision
of MSD.

% Native data; generator: figures/source/dataset-methods-examples.py.
% Provenance: reports/thesis_dataset_methods_examples_20261103/provenance.json.
\begin{figure}[htbp]
\centering
\includegraphics[width=0.98\linewidth]{figures/compiled/adni-data-examples.pdf}
\caption{Paired ADNI examples from one subject: the left and right
hippocampal crops, each shown as an MRI slice, a reference mask and an
annotation overlay. Blue denotes hippocampus on both sides; white in the
masks denotes background. Each row is a slice through a separate
$32\times48\times40$ volume. MRI contrast is adjusted separately for each
crop using its first and ninety-ninth intensity percentiles. The panels
follow the stored grids; no anatomical plane or physical scale is inferred
from the exported coordinate metadata. These are reference annotations,
not model predictions.}
\label{fig:adni-data-examples}
\end{figure}

Crops are centred on the reference-mask bounding box, with identical
indices applied to the image. Their evaluation therefore assumes prior
hippocampal localization. The capped experiments use three partitions,
identified as folds 0, 1 and 2. Each partition groups both sides of a
subject together, with an eligible training pool of 108 subjects
(216 crops) and 27 validation subjects (54 crops), with no within-fold
subject overlap. Five selected training subjects contribute both sides:
ten support crops per fold, shared across models and seeds 0, 1 and 4.
Images are
z-score normalized over nonzero intensities and padded to $64^3$,
without resizing. MedSAM3 uses the separate
two-crop support protocol in Section~\ref{sec:medsam-adni-methods}.

The exported ADNI crops do not retain the original spacing and orientation
metadata, so spatial constraints are defined in voxel coordinates. They
retain the same geometric construction used for MSD: signed distances are
capped at eight voxels, and depth supervision uses layers two voxels thick
within the hippocampus and its eight-voxel surrounding region. The depth
plane passes through the hippocampus centre and contains its main axis; in
ADNI, an image-grid axis fixes its orientation. Because MSD has verified
isotropic spacing, meaning the same $1\,\mathrm{mm}$ spacing along all three
axes, its $2\,\mathrm{mm}$ layers and $8\,\mathrm{mm}$ surrounding region
match these widths numerically on the voxel grid. The two formulations are
therefore comparable in construction and grid-scale settings, without
assuming identical physical spacing or anatomical orientation across the
datasets.




\subsection{Four-Model Comparisons under a FLOP Ceiling}
\label{sec:main-msd-methods}

\subsubsection{Models and Baseline Objectives}

nnU-Net, SwinUNETR, UNETR and a compact three-dimensional U-Net start
from random initialization, use batch size one and predict the same
dataset-specific $64^3$ grid: three classes on MSD and two on ADNI. Their configurations instantiate the model
families introduced in Section~\ref{sec:segmentation-architectures}.
MSD parameter counts and shared optimizer settings are listed in
Table~\ref{tab:main-model-settings}.

This selection compares a compact convolutional network, a configured
convolutional segmentation system and two ways of incorporating attention.
Keeping the same output grid and labelled cases makes their responses
to the added supervision comparable within the stated compute allowance.
Layer widths and other implementation settings are recorded in
Appendix~\ref{sec:conventional-implementation}.

The native baseline objectives are retained
(Table~\ref{tab:capped-baseline-losses}). The Medical Open Network for
Artificial Intelligence (MONAI) Dice loss uses all dataset
softmax classes with a linear denominator. nnU-Net combines all-class
cross-entropy and negative mean foreground Dice, pooling region sums
over the batch. Both implementations use $10^{-5}$ numerator and
denominator smoothing. Negative Dice and one minus Dice differ by a
constant and produce the same derivative.

\begin{table}[htbp]
\centering
\small
\begin{tabular}{P{4.3cm}P{3.2cm}P{6.0cm}}
\hline
Model & Baseline objective & Training Dice classes and supervision \\
\hline
SwinUNETR, UNETR and 3D U-Net
& Soft Dice
& All dataset classes, including background; final output. \\[0.5em]
nnU-Net
& Soft Dice plus cross-entropy
& Foreground classes in Dice; all classes in cross-entropy;
  multiple output resolutions. \\
\hline
\end{tabular}
\caption{Baseline losses in the capped MSD and ADNI comparisons. The class set
listed for Dice concerns training, not the selection of evaluation metrics.}
\label{tab:capped-baseline-losses}
\end{table}

nnU-Net also retains its weighted losses at multiple output resolutions,
with nearest-neighbour resampling of labels. Auxiliary terms act on
the primary, full-resolution prediction. The paired comparison within
each model thus isolates a change to its configured training objective;
cross-model comparisons also reflect different baselines and optimizers.

Each model follows a fixed learning-rate trajectory across its
objectives. If an extra loss reduces the number of updates affordable
under the cap, the schedule is not shortened to compensate. The
comparison therefore measures the accuracy reached with the complete
objective inside the allowance. Optimizer schedules and precision
settings are given in Appendix~\ref{sec:conventional-implementation}.

Fourteen objectives cover the combinations in
Table~\ref{tab:constraint-combinations}; $\varnothing$ denotes no
selected term among $B,E,S$. Bands and edges supervise
the hippocampal union on MSD and the single foreground on ADNI; SDF
supervises each foreground class, and depth
uses the original reference classes. Preliminary MSD screening did not
establish a consistent incremental benefit from edge alone over the
supervised baseline, or from edge plus SDF over SDF alone. These two
configurations were therefore omitted from the subsequent
fourteen-objective replications across datasets. Their exclusion is a
scope decision based on the preliminary comparisons; the main replication
provides evidence for the retained configurations. It consequently
supports comparisons among those configurations, rather than a test of every possible combination and interaction.
% EXCLUSION EVIDENCE: reports/compute_matched_constraints_20pct_20261004/
% summary.json (U-Net, MSD fold 0 / seed 0): E 41.33 vs baseline 42.14;
% E+S 57.25 vs S 61.15, although E+S exceeds baseline. The longer pilot in
% reports/compute_matched_constraints_682338/README.md records only a
% +0.055 pp validation-selected edge gain and a slightly lower final score.
% These archives do not establish ineffectiveness across all datasets/seeds.
% Pilot records are preserved; omitted configurations were not part of the
% later fourteen-objective replication inventory.

\begin{table}[htbp]
\centering
\small
\setstretch{1.05}
\renewcommand{\arraystretch}{1.2}
\begin{tabular}{@{}P{4.6cm}P{4.6cm}P{4.6cm}@{}}
\hline
Subset of $B,E,S$ & Without depth & With depth \\
\hline
$\varnothing$ & Baseline & Baseline $+D$ \\
$B$ & Baseline $+B$ & Baseline $+B+D$ \\
$E$ & Not replicated & Baseline $+E+D$ \\
$S$ & Baseline $+S$ & Baseline $+S+D$ \\
$B,E$ & Baseline $+B+E$ & Baseline $+B+E+D$ \\
$B,S$ & Baseline $+B+S$ & Baseline $+B+S+D$ \\
$E,S$ & Not replicated & Baseline $+E+S+D$ \\
$B,E,S$ & Baseline $+B+E+S$ & Baseline $+B+E+S+D$ \\
\hline
\end{tabular}
\caption{The fourteen objective combinations in each conventional dataset
comparison and in the SAM3 and MedSAM3 adaptation campaigns.
Each included entry retains the architecture's supervised
baseline. Coefficients are calibrated for the corresponding combination;
adding a family need not leave every existing coefficient unchanged.}
\label{tab:constraint-combinations}
\end{table}

\subsubsection{Calibration and Activation}
\label{sec:calibration-methods}

An auxiliary loss can dominate an update simply because its numerical
scale differs from the baseline. Calibration sets its coefficient by
comparing the size of their gradients on the labelled support cases.
This uses training annotations only and holds the measured network state
fixed. The coefficients are then retained during optimization.

The medical studies calibrate after a short baseline prefix: nine updates
for nnU-Net and 50 for the other models. In a combined objective, terms
are calibrated sequentially against the baseline and preceding terms.
The order is bands, edges, signed distance and depth. Consequently, a
term's coefficient can change when it is used in a different combination;
these experiments compare complete calibrated objectives. A median
gradient target sets the typical strength, while an upper-percentile
bound limits unusually large contributions. The equations and numerical
settings are given in Appendix~\ref{sec:detailed-calibration}.

Auxiliary terms enter gradually after the prefix. Signed distance either
shares that ramp or starts later, according to the fixed activation
policy in Table~\ref{tab:sdf-activation-policies}. This policy was developed
from earlier MSD validation experiments and retained for the reported
campaigns; it was not selected independently of those exploratory results.
Combinations containing both signed distance and depth use delayed
activation of the distance term.

\subsubsection{Resource Accounting and Evaluation}

The medical ceiling corresponds to the measured cost of 90 nnU-Net
updates, approximately 31,976 billion floating-point operations (GFLOPs).
Exploratory learning curves informed this short reference allowance.
With ten support crops and one crop per update, the reference represents
nine passes through the training subset. nnU-Net has the most expensive
update among the four implementations, so the same allowance lets models
with cheaper updates perform more optimization steps.

Every medical model and objective receives that ceiling. The accounting
includes the forward pass, active training losses, backpropagation and
the optimizer update. A step is allowed only when its full cost fits in
the remaining budget. Reference-target construction, validation, data
loading and storage are excluded, and calibration is recorded separately.
This matches optimization work, not elapsed time or end-to-end computation.
Equal allowances also need not place the architectures equally close to
convergence.
Figure~\ref{fig:flop-budget-example} illustrates how update cost affects
training exposure and distinguishes the medical and geological
calibration accounting.

\begin{figure}[htbp]
\centering
% Generator: figures/source/explanatory-concepts.py
\includegraphics[width=\linewidth]{figures/compiled/flop-budget-example.pdf}
\caption{Computation budgets in the comparisons from scratch.
(a) Each block bounded by white separators represents one complete
optimizer update. A more expensive update permits fewer steps under
the same ceiling; auxiliary computation also consumes that allowance.
Costs and update counts here are illustrative, not experimental
measurements. (b) Medical comparisons report calibration separately
from their optimization ceiling. The geological ceiling includes both
calibration and optimization. The outlined containers indicate
accounting scope, not relative measured costs or equal budgets across
datasets.}
\label{fig:flop-budget-example}
\end{figure}
\FloatBarrier

The reported checkpoint is the final state admitted by the budget, with
no early stopping or selection of a validation-best checkpoint. Selecting
the highest-mean objective after the experiments is a separate, descriptive
comparison. All 504 runs in each medical study enter the aggregation.
Appendix~\ref{sec:detailed-resource-accounting} gives the exact ceiling
and arithmetic conventions; Table~\ref{tab:cap-training-exposure} reports
how the resulting update counts translate into passes through the data.
The same numerical medical ceiling is used on MSD and ADNI, while their
different label vocabularies remain part of the interpretation.

\subsection{Adapting Pretrained Models}

The adaptation studies ask how much the same spatial supervision helps
when an image representation is already available. All three studies
freeze the SAM3 base weights and train only LoRA factors, giving
18,497,664 trainable adapter parameters in each study. The geological
study starts from fresh factors with a zero initial update; the medical
studies start from the released MedSAM3 adapter
\cite{carion2025sam3,liu2025medsam3}. Within each matched comparison, the
same base weights and initial adapter parameters are restored before
every baseline and constrained run. These runs also share the support
images and sampling schedule, so their comparison measures the effect
of changing the training objective from the same starting point.

The pretrained baseline supervises several outputs. Classification and
bounding-box losses assess detected objects and their enclosing boxes; a
presence loss assesses whether the queried target occurs in the image. The
mask losses combine soft Dice with focal loss, a form of cross-entropy that
reduces the weight of already confident correct predictions. Spatial
penalties are added to that objective. Their coefficients are calibrated
from gradients on the support annotations. Each optimizer update uses two
adjacent-slice pairs, allowing edge supervision to compare predictions
across planes as well as within them. Reference supports are constructed in
three dimensions. Full-volume counts and sampling corrections ensure that
averaging over all possible sampled planes recovers the corresponding volume
loss. Normalizing by each sampled plane's own support count would give a
different weighting. Appendix~\ref{sec:sampled-pretrained-supervision} gives
these reductions.

Inference uses image planes and fixed text queries. The model's confidence
in each mask location, its match to the text query and the target's presence
supply the foreground responses, which are converted to mutually exclusive
labels and reassembled into volumes. Reference masks are used to score these
predictions, never to provide localization prompts. Each study evaluates the
adapter at its final update, without early stopping or validation-based
checkpoint selection. Dice is averaged within a run, then across seeds
within each fold and equally across folds.

\paragraph{SAM3 on GeoFWI3D.}\label{sec:sam3-methods}
Each fold uses five labelled support volumes, taken from the beginning
of its established support list, and 96 held-out volumes. The native
$96^3$ reflectivity volumes are mapped from $[-1,1]$ to $[0,1]$ for
image processing. The query \emph{salt body} produces one foreground
response; its complement represents background.

Every objective receives 25 updates, using the first 25 steps of the
retained 200-update learning-rate schedule. The schedule is not compressed
into a shorter decay. Bands are calibrated at initialization; other selected
terms are calibrated after five prefix updates, during which their weights
are zero. Bands and edges use the salt mask, with two-step, zero-padded
morphology retaining crop-face support. SDF uses one foreground field on the
unit voxel grid. Depth weights use source-array geometry, with no anatomical
orientation assumption. These choices distinguish this adaptation protocol
from the conventional geological recipes.

\paragraph{MedSAM3 on MSD.}\label{sec:medsam-methods}
Five support crops are used in each fold, with 46, 52 and 44 held-out
crops respectively. The 15 distinct support crops and 142 distinct
evaluation crops do not overlap. Fold selection followed earlier
exploratory outcomes, and participant identities are unavailable; the
crop separation therefore does not establish participant independence.
Images and masks are reordered and, where needed, flipped together
into the closest right--anterior--superior orientation recorded by their
NIfTI coordinate transforms. Intensities are clipped and scaled between the first and
99th percentiles, then converted to eight-bit planes.

The fixed queries are \emph{anterior hippocampus} and \emph{posterior
hippocampus}. Each response probability $q$ is converted to log-odds,
$\log(q/(1-q))$. These scores, together with a zero background logit, are
normalized by softmax before argmax assigns one class per voxel. Bands and
edges are applied separately to background, anterior and posterior, so they
include the internal interface. SDF uses the two foreground classes and
depth uses physical geometry. All variants receive 200 updates, with their
selected auxiliary terms ramping from the first update.

The fourteen objectives were completed across two execution batches.
The original bands and bands-plus-edge results retain their matched
baseline and calibration; the other eleven constrained variants use
the controls and calibration recorded with their batch. Those controls
are not pooled, and repeated baseline runs do not create additional
objective configurations.
Figure~\ref{fig:medsam3-workflow} shows the adaptation and prediction
path.

\begin{figure}[htbp]
\centering
\setstretch{1.05}
\begin{tikzpicture}[x=1cm,y=-1cm,font=\footnotesize,
  every node/.style={text=thesisInk},>=stealth]
\node[align=center,text width=15cm,inner sep=0pt,text=thesisGray,font=\fontsize{9}{11}\selectfont\bfseries] at (7.750,0.200) {TASK-SPECIFIC ADAPTATION ON MSD};
\draw[fill=thesisPaleGray,draw=thesisGray,line width=.65pt,rounded corners=2pt] (0.150,0.625) rectangle (7.450,1.775);
\node[align=center,text width=7.06cm,inner sep=0pt,text=thesisInk,font=\fontsize{10}{12}\selectfont] at (3.800,1.200) {SAM3 base weights\\Frozen throughout adaptation};
\draw[fill=thesisPaleOrange,draw=thesisOrange,line width=.65pt,rounded corners=2pt] (8.050,0.625) rectangle (15.350,1.775);
\node[align=center,text width=7.06cm,inner sep=0pt,text=thesisInk,font=\fontsize{10}{12}\selectfont] at (11.700,1.200) {Released medical LoRA adapter\\Starting factors restored for each arm};
\draw[draw=thesisGray,line width=0.8pt,->] (3.800,1.775) -- (3.800,2.150) -- (7.750,2.150) -- (7.750,2.500);
\draw[draw=thesisGray,line width=0.8pt] (11.700,1.775) -- (11.700,2.150) -- (7.750,2.150);
\draw[fill=thesisPaleBlue,draw=thesisBlue,line width=.65pt,rounded corners=2pt] (0.150,2.500) rectangle (15.350,3.600);
\node[align=center,text width=14.959999999999999cm,inner sep=0pt,text=thesisInk,font=\fontsize{10}{12}\selectfont] at (7.750,3.050) {MedSAM3 initialization\\Pretrained base + medical adapter};
\draw[draw=thesisGray,line width=0.8pt,->] (11.700,3.600) -- (11.700,4.150);
\draw[fill=thesisPaleBlue,draw=thesisBlue,line width=.65pt,rounded corners=2pt] (0.150,4.150) rectangle (7.150,6.250);
\node[align=center,text width=6.76cm,inner sep=0pt,text=thesisInk,font=\fontsize{10}{12}\selectfont] at (3.650,5.200) {5 labelled support volumes\\2 adjacent-slice pairs per update\\Fixed anterior / posterior queries};
\draw[fill=thesisPaleOrange,draw=thesisOrange,line width=.65pt,rounded corners=2pt] (8.050,4.150) rectangle (15.350,6.250);
\node[align=center,text width=7.06cm,inner sep=0pt,text=thesisInk,font=\fontsize{10}{12}\selectfont] at (11.700,5.200) {Baseline / baseline + constraints\\200 updates; only LoRA factors (18.50 M)\\14 objectives; two execution batches};
\draw[draw=thesisGray,line width=0.8pt,->] (7.150,5.200) -- (8.050,5.200);
\draw[draw=thesisGray,line width=0.3pt] (0.100,6.700) -- (15.400,6.700);
\node[align=center,text width=9.5cm,inner sep=0pt,text=thesisGray,font=\fontsize{9}{11}\selectfont\bfseries] at (5.000,7.150) {INFERENCE ON HELD-OUT VOLUMES};
\draw[draw=thesisGray,line width=0.8pt,->,dashed] (11.700,6.250) -- (11.700,7.800);
\draw[fill=thesisPaleBlue,draw=thesisBlue,line width=.65pt,rounded corners=2pt] (0.150,7.800) rectangle (7.150,9.100);
\node[align=center,text width=6.76cm,inner sep=0pt,text=thesisInk,font=\fontsize{10}{12}\selectfont] at (3.650,8.450) {MRI slices\\Fixed anterior / posterior queries};
\draw[fill=thesisPaleBlue,draw=thesisBlue,line width=.65pt,rounded corners=2pt] (8.050,7.800) rectangle (15.350,9.100);
\node[align=center,text width=7.06cm,inner sep=0pt,text=thesisInk,font=\fontsize{10}{12}\selectfont] at (11.700,8.450) {Adapted MedSAM3\\Two foreground response maps};
\draw[draw=thesisGray,line width=0.8pt,->] (7.150,8.450) -- (8.050,8.450);
\draw[draw=thesisGray,line width=0.8pt,->] (11.700,9.100) -- (11.700,9.800);
\draw[fill=thesisPaleBlue,draw=thesisBlue,line width=.65pt,rounded corners=2pt] (8.050,9.800) rectangle (15.350,11.400);
\node[align=center,text width=7.06cm,inner sep=0pt,text=thesisInk,font=\fontsize{10}{12}\selectfont] at (11.700,10.600) {Three-class semantic projection\\Background + anterior + posterior\\Voxelwise argmax};
\draw[fill=thesisPaleGray,draw=thesisGray,line width=.65pt,rounded corners=2pt] (0.150,9.800) rectangle (7.150,11.400);
\node[align=center,text width=6.76cm,inner sep=0pt,text=thesisInk,font=\fontsize{10}{12}\selectfont] at (3.650,10.600) {Reassembled 3D label map\\0: background; 1: anterior; 2: posterior};
\draw[draw=thesisGray,line width=0.8pt,->] (8.050,10.600) -- (7.150,10.600);
\node[align=center,text width=15.3cm,inner sep=0pt,text=thesisGray,font=\fontsize{9}{11}\selectfont] at (7.750,11.850) {Evaluation: Dice of final predictions; paired gains compare each variant with its baseline.};
\end{tikzpicture}
\caption{MedSAM3 adaptation and inference on MSD. Training changes only
the medical LoRA factors using five support crops per fold. The two
fixed queries produce foreground responses, which are combined with
background to obtain mutually exclusive labels. Held-out masks are used
for evaluation and supply no localization prompts.}
\label{fig:medsam3-workflow}
\end{figure}

\paragraph{MedSAM3 on ADNI.}\label{sec:medsam-adni-methods}
Each fold uses two labelled crops from two subjects and 54 held-out
crops from 27 other subjects. Folds 0 and 3 share their two support
crops; fold 1 shares one and changes the other. Only three distinct
support crops occur across the study. These folds vary the evaluation
partition, but do not provide three independent draws of a training set.

The medical initialization, 200-update schedule and intensity scaling
follow the MSD protocol. ADNI retains its stored array orientation and
uses the single query \emph{hippocampus}. The binary projection combines
the foreground response with its complement. Bands and edges supervise
the hippocampus mask, without a second background-mask average.
SDF uses unit array voxels, and depth uses source-axis geometry without
an inferred anatomical orientation.

These update budgets describe additional task adaptation and exclude
pretraining. They do not provide a FLOP match with the conventional
studies. Exact optimizer settings, calibration rules, fold and seed
identifiers, and execution environments are retained in
Appendix~\ref{sec:adaptation-implementation}. Runtime comparisons keep
their recorded scope; a training-loop timer is not an end-to-end
adaptation cost.

\subsection{Evaluation Measures}
\label{sec:evaluation-methods}

Dice serves two distinct roles. During training, a differentiable soft
Dice loss contributes to the architecture-specific optimization
objective. During evaluation, hard Dice measures overlap between the
final discrete predictions and reference labels. Every Dice score in
the results figures and tables uses this second definition, including scores on the
training support set. Validation Dice uses held-out examples;
training-set Dice uses examples supplied during fitting. Neither score
contains focal, cross-entropy, box, presence or spatial-constraint loss
terms, and neither is the value of the Dice loss used during training.

Voxelwise argmax defines the predicted label. For reference region
$Y_{i,c}$ and prediction $\widehat Y_{i,c}$, hard class Dice is
\begin{equation}
 D_{i,c}=\frac{2|Y_{i,c}\cap\widehat Y_{i,c}|}
 {|Y_{i,c}|+|\widehat Y_{i,c}|}.
\label{eq:reported-hard-dice}
\end{equation}
A present reference class with no predicted voxels scores zero.
The MSD volume score is the mean anterior/posterior Dice,
$D_i=(D_{i,A}+D_{i,P})/2$, and the ADNI volume score is hippocampus Dice. Geological evaluation
uses salt Dice.
Background is excluded from these primary metrics. Volumes receive
equal weight in a run mean. Section~\ref{sec:background-dice-results}
adds background and all-class Dice as complementary measures on
the same predictions, retaining this order of averaging.

Let $\overline D_{m,j,f,s}$ denote the mean volume Dice in a run with
model $m$, objective $j$, fold $f$ and seed $s$. Subscript
$\mathrm{base}$ identifies its matched baseline. Their paired change is
\begin{equation}
 \Delta_{m,j,f,s}=100\left(
 \overline D_{m,j,f,s}-\overline D_{m,\mathrm{base},f,s}\right).
\label{eq:paired-reported-gain}
\end{equation}
Dice is reported on a percentage scale and differences in percentage
points. Reported standard deviations are sample standard deviations across
the three seed-averaged fold means; paired-gain dispersions use the three
fold-mean differences. Positive-change counts use all nine matched
fold/seed cells. Each figure or table identifies its reduction explicitly.
These standard deviations are not confidence intervals over independent datasets: seeds share data,
training pools overlap and configuration selection is exploratory.
Dispersion across cases is identified separately from dispersion
across fold/seed cells.

A confusion matrix counts how often each reference class is assigned each
predicted class. The pretrained-model matrices use reference classes as rows
and predicted classes as columns. For MSD, four disjoint errors are counted:
background predicted as foreground, foreground predicted as background,
anterior predicted as posterior, and posterior predicted as anterior. Their
sum gives all incorrectly labelled voxels. Conventional-run error counts
pool the full validation grid before averaging runs. MedSAM3 tables average
errors per crop within each run and then weight folds equally. All counts
cover the complete prediction, including voxels outside the boundary bands.
SAM3 FN/FP are likewise averaged per held-out volume within a run, then over
seeds and folds; each run contains 96 validation volumes. Binary ADNI and
GeoFWI3D have no A/P swaps. The different error-count reductions are
labelled and are not compared as if they had the same denominator.

\clearpage

